Partition \([0,1]^d\) into \(K\le(2/\delta)^d\) half-open cubes \(B_1,\ldots,B_K\) of diameter at most \(C_d\delta\), and split the sample into folds of sizes \(N_1,N_2\asymp n\). On the first fold let \(\widehat p_j=N_1^{-1}\sum_{i\in I_1}\mathbf1\{X_i\in B_j\}\). On the second, let \(\widehat m_j\) average the outcomes with \(X_i\in B_j\) and \(\lvert A_i-t_0\rvert\le h\), assigning zero to empty cells. Define the Borel estimator
\[
\widehat\theta_{h,\delta}=\sum_{j=1}^K\widehat p_j\widehat m_j.
\]
Write \(p_j=P_X(B_j)\), \(q_h(x)=P(\lvert A-t_0\rvert\le h\mid X=x)\),
\[
m_j=\mathbb E[Y\mid X\in B_j,\lvert A-t_0\rvert\le h],
\qquad
\bar\mu_j=\mathbb E[\mu_P(t_0,X)\mid X\in B_j].
\]
Positivity gives \(2ch\le q_h(x)\le2h/c\). Treatment Hölder continuity gives
\[
\left|q_h(x)^{-1}\int_{t_0-h}^{t_0+h}
\mu_P(a,x)\pi_P(a\mid x)\,da-\mu_P(t_0,x)\right|
\le Ch^{\bar\alpha},
\]
and
\[
q_h(x)/(2h)=\pi_P(t_0\mid x)+r_h(x),
\qquad \|r_h\|_\infty\le Ch^{\bar\beta}.
\]
Within a cell, both fixed-slice functions oscillate by at most \(C\delta^{\bar s}\). The difference between a \(q_h\)-weighted cell mean and an unweighted one is their covariance divided by the cell mean of \(q_h\). Decomposing \(q_h/(2h)\) as above and bounding the covariance by the product of the two oscillations yields
\[
|m_j-\bar\mu_j|
\le C\{h^{\bar\alpha}+\delta^{2\bar s}+\delta^{\bar s}h^{\bar\beta}\}.
\]
Since \(\theta_P(t_0)=\sum_jp_j\bar\mu_j\), the squared approximation bias is at most
\[
C\{h^{2\bar\alpha}+\delta^{4\bar s}+\delta^{2\bar s}h^{2\bar\beta}\}.
\]
Conditional on the second fold, the first-fold contribution has variance at most \(M_Y^2/N_1\). Put \(Q_j=P(X\in B_j,|A-t_0|\le h)\ge2chp_j\). If \(N_j\sim\operatorname{Bin}(N_2,Q_j)\) is the selected count, then
\[
\mathbb E\!\left[\frac{\mathbf1\{N_j>0\}}{N_j}\right]
\le\frac{2}{N_2Q_j}.
\]
Thus nonempty-cell variance is at most \(C\sum_jp_j^2/(N_2Q_j)\le C/(nh)\). For the empty-cell mass \(W=\sum_jp_j\mathbf1\{N_j=0\}\),
\[
\mathbb E[W^2]\le\mathbb E[W]
\le\sum_jp_je^{-N_2Q_j}
\le\frac{CK}{nh}
\le\frac{C}{nh\delta^d}.
\]
Combining these bounds proves the risk inequality. Substitution of the two power tunings and selection of the smallest exponent gives \(R(u,v)\).